\lyseis{Α3.6}{Δυνάμεις ρητών με εκθέτη μη μηδενικό φυσικό}
\sol{1} α) Σωστό. β) Λάθος: $\frac{2^3}{3^3}=\frac{8}{27}$. γ) Λάθος: $0{,}2^4=0{,}0016$ ενώ $4\cdot0{,}2=0{,}8$. δ) Σωστό: $\left(-\frac{3}{10}\right)^2=\frac{9}{100}=0{,}09=0{,}3^2$. ε) Σωστό: το πρόσημο είναι $-$ (περιττός εκθέτης) και $\left(\frac{1}{3}\right)^3=\frac{1}{27}$.
\sol{2} α) $\frac{1}{4}$ \quad β) $-\frac{1}{27}$ \quad γ) $\frac{4}{25}$ \quad δ) $0{,}04$ \quad ε) $0{,}027$ \quad στ) $-1$.
\sol{3} α) Αρνητικό (αρνητική βάση, περιττός εκθέτης). β) Θετικό (άρτιος). γ) Θετικό (άρτιος). δ) Αρνητικό: η δύναμη $\left(\frac{3}{5}\right)^3$ είναι θετική και το $-$ μπροστά την κάνει αρνητική. ε) Αρνητικό (περιττός).
\sol{4} α) $\left(\frac{3}{5}\right)^3$ \quad β) $(-2)^4$ \quad γ) $0{,}1^2$ \quad δ) $\left(-\frac{1}{2}\right)^2$.
\sol{5} α) $(-2)^4=16$ και $-2^4=-16$: διαφέρουν. β) $(-3)^3=-27$ και $-3^3=-27$: ίσα. γ) $\left(-\frac{1}{2}\right)^2=\frac{1}{4}$ και $-\left(\frac{1}{2}\right)^2=-\frac{1}{4}$: αντίθετα.
\sol{6} α) $1$ \quad β) $0$ \quad γ) $1$ \quad δ) $-1$ \quad ε) $1\,000$ \quad στ) $0{,}01$.
\sol{7} α) $\frac{1}{4}-\frac{1}{8}+\frac{1}{16}=\frac{4}{16}-\frac{2}{16}+\frac{1}{16}=\frac{3}{16}$. β) $0{,}01+0{,}01-0{,}04=-0{,}02$. γ) $-1+1+1-1=0$.
\sol{8} Μετά την τρίτη: $2\cdot\left(\frac{1}{2}\right)^3=\frac{1}{4}$ μ. Μετά την τέταρτη: $2\cdot\left(\frac{1}{2}\right)^4=2\cdot\frac{1}{16}=\frac{1}{8}$ μ.
\sol{9} $\left(\frac{1}{2}\right)^5=\frac{1}{2^5}=\frac{1}{32}$ του αρχικού εμβαδού.
\sol{10} $1{,}5^2=1{,}5\cdot1{,}5=2{,}25$ τετραγωνικά μέτρα. Ο Ηλίας πολλαπλασίασε τη βάση με τον εκθέτη αντί να τη χρησιμοποιήσει ως παράγοντα δύο φορές.
\sol{11} $8\cdot\left(-\frac{1}{8}\right)=-1$ και $(0{,}8)^2=0{,}64$. Άθροισμα: $-1+0{,}64=-0{,}36$.
\sol{12} $0{,}1$, \ $0{,}01$, \ $0{,}001$, \ $0{,}0001$. Κάθε δύναμη έχει τόσα δεκαδικά ψηφία όσος είναι ο εκθέτης, με το $1$ στην τελευταία θέση.
\sol{13} Το $-3^2$ είναι ο αντίθετος του $3^2$, δηλαδή $-9$. Για να υψώσουμε τον $-3$, χρειάζεται παρένθεση: $(-3)^2=(-3)(-3)=+9$, γιατί δύο αρνητικοί δίνουν θετικό.
\sol{14} Όταν ο εκθέτης είναι άρτιος. Οι αρνητικοί παράγοντες μπορούν να χωριστούν σε ζεύγη και κάθε ζεύγος δίνει θετικό γινόμενο.
\sol{15} $\left(\frac{1}{2}\right)^3=\frac{1}{8}<\frac{1}{4}=\left(\frac{1}{2}\right)^2$. Στη βάση $2$: $2^3=8>4=2^2$. Με βάση μικρότερη του $1$ οι δυνάμεις μικραίνουν όσο μεγαλώνει ο εκθέτης, ενώ με βάση μεγαλύτερη του $1$ μεγαλώνουν.
\sol{16} Το $\alpha^2$ είναι θετικό (άρτιος εκθέτης) και το $\alpha^3$ αρνητικό (αρνητική βάση, περιττός εκθέτης), άρα $\alpha^2>\alpha^3$.
\sol{17} $\frac{1}{2}=0{,}5$, \ $\frac{1}{4}=0{,}25$, \ $\frac{1}{8}=0{,}125$, \ $\frac{1}{16}=0{,}0625$, \ $\frac{1}{32}=0{,}03125$, \ $\frac{1}{64}=0{,}015625$. Οι δυνάμεις μικραίνουν, και κάθε μία είναι το μισό της προηγούμενης.
\sol{18} Ελεύθερη απάντηση. Ο μαθητής πρέπει να τηρεί την προτεραιότητα (παρενθέσεις, δυνάμεις, πολλαπλασιασμοί, προσθέσεις) και να ελέγχει το πρόσημο κάθε δύναμης.
